\begin{lemma}
\label{lemma-rhe}
Notation and assumptions as in Lemma \ref{lemma-rh}. For a closed point
$x \in X$ let $d_x$ be the multiplicity of $x$ in $R$. Then
$$
2g_X - 2 = (2g_Y - 2) \deg(f) + \sum\nolimits d_x [\kappa(x) : k]
$$
Moreover, we have the following results
\begin{enumerate}
\item $d_x = \text{length}_{\mathcal{O}_{X, x}}(\Omega_{X/Y, x})$,
\item $d_x \geq e_x - 1$ where $e_x$ is the ramification index
of $\mathcal{O}_{X, x}$ over $\mathcal{O}_{Y, y}$,
\item $d_x = e_x - 1$ if and only if $\mathcal{O}_{X, x}$ is tamely
ramified over $\mathcal{O}_{Y, y}$.
\end{enumerate}
\end{lemma}

\begin{proof}
By Lemma \ref{lemma-rh} and the discussion above
(which used
Varieties, Lemma \ref{varieties-lemma-algebraic-scheme-dim-0}
and
Algebra, Lemma \ref{algebra-lemma-pushdown-module})
it suffices to prove the results on the
multiplicity $d_x$ of $x$ in $R$. Part (1) was proved
in the discussion above. In the discussion above
we proved (2) and (3) only in the case where $\kappa(x)$ is
separable over $k$.
In the rest of the proof we give a uniform treatment
of (2) and (3) using material on differents of
quasi-finite Gorenstein morphisms.

\medskip\noindent
First, observe that $f$ is a quasi-finite Gorenstein morphism.
This is true for example because
$f$ is a flat quasi-finite morphism and $X$ is Gorenstein
(see Duality for Schemes, Lemma
\ref{duality-lemma-flat-morphism-from-gorenstein-scheme})
or because it was shown in the proof of
Discriminants, Lemma
\ref{discriminant-lemma-discriminant-quasi-finite-morphism-smooth}
(which we used above). Thus $\omega_{X/Y}$ is invertible by
Discriminants, Lemma \ref{discriminant-lemma-gorenstein-quasi-finite}
and the same remains true after replacing $X$ by opens and after
performing a base change by some $Y' \to Y$. We will use this
below without further mention.

\medskip\noindent
Choose affine opens $U \subset X$ and $V \subset Y$
such that $x \in U$, $y \in V$, $f(U) \subset V$, and $x$ is the only
point of $U$ lying over $y$. Write $U = \Spec(A)$ and $V = \Spec(B)$.
Then $R \cap U$ is the different of $f|_U : U \to V$.
By Discriminants, Lemma \ref{discriminant-lemma-base-change-different}
formation of the different commutes with arbitrary base change
in our case. By our choice of $U$ and $V$ we have
$$
A \otimes_B \kappa(y) =
\mathcal{O}_{X, x} \otimes_{\mathcal{O}_{Y, y}} \kappa(y) =
\mathcal{O}_{X, x}/(s^{e_x})
$$
where $e_x$ is the ramification index as in the statement of the lemma.
Let $C = \mathcal{O}_{X, x}/(s^{e_x})$ viewed as a finite algebra
over $\kappa(y)$. Let $\mathfrak{D}_{C/\kappa(y)}$ be the different
of $C$ over $\kappa(y)$ in the sense of
Discriminants, Definition \ref{discriminant-definition-different}.
It suffices to show: $\mathfrak{D}_{C/\kappa(y)}$
is nonzero if and only if the extension
$\mathcal{O}_{Y, y} \subset \mathcal{O}_{X, x}$ is tamely ramified
and in the tamely ramified case $\mathfrak{D}_{C/\kappa(y)}$
is equal to the ideal generated by $s^{e_x - 1}$ in $C$.
Recall that tame ramification means exactly that $\kappa(x)/\kappa(y)$
is separable and that the characteristic of $\kappa(y)$ does not
divide $e_x$. On the other hand, the different of $C/\kappa(y)$ is nonzero
if and only if $\tau_{C/\kappa(y)} \in \omega_{C/\kappa(y)}$ is nonzero.
Namely, since $\omega_{C/\kappa(y)}$ is an invertible $C$-module
(as the base change of $\omega_{A/B}$)
it is free of rank $1$, say with generator $\lambda$. Write
$\tau_{C/\kappa(y)} = h\lambda$ for some $h \in C$. Then
$\mathfrak{D}_{C/\kappa(y)} = (h) \subset C$ whence the claim.
By Discriminants, Lemma \ref{discriminant-lemma-tau-nonzero}
we have $\tau_{C/\kappa(y)} \not = 0$
if and only if $\kappa(x)/\kappa(y)$
is separable and $e_x$ is prime to the characteristic.
Finally, even if $\tau_{C/\kappa(y)}$ is nonzero, then
it is still the case that $s \tau_{C/\kappa(y)} = 0$
because $s\tau_{C/\kappa(y)} : C \to \kappa(y)$
sends $c$ to the trace of the nilpotent operator $sc$ which is zero.
Hence $sh = 0$, hence $h \in (s^{e_x - 1})$ which proves
that $\mathfrak{D}_{C/\kappa(y)} \subset (s^{e_x - 1})$ always.
Since $(s^{e_x - 1}) \subset C$ is the smallest nonzero ideal,
we have proved the final assertion.
\end{proof}